% Source MD: E:\wanru\PRML_Course\NoteMD\Week1\第4节课 模型优化中的局部最小与鞍点问题第4节课 模型优化中的局部最小与鞍点问题.md
\documentclass[main.tex]{subfiles}

\graphicspath{{\subfix{figures/}}}

\begin{document}

\chapter{模型优化中的局部最小与鞍点问题}
\label{chap:w01l04-local-minima-and-saddle-points}

\begin{warningbox}
第四节课的原始 Markdown 文件名中存在重复标题；TeX 工程侧已经改用稳定 slug，
后续迁移正文时无需继承这个命名问题。
\end{warningbox}

\section{临界点分析骨架}
\label{sec:w01l04-critical-points}

\begin{theorem}{Hessian 判别思路}{w01l04-hessian}
如果某个临界点处的 Hessian 矩阵特征值有正有负，那么该点是鞍点而不是局部最小值。
\end{theorem}

\begin{equation}
  L\paren{\theta} \approx L\paren{\theta'}
  + \frac{1}{2}(\theta - \theta')\trans
  H
  (\theta - \theta').
  \label{eq:w01l04-taylor}
\end{equation}

\begin{lstlisting}[style=text,caption={课堂整理占位模板},label={lst:w01l04-outline}]
1. 先写问题背景
2. 再写 Taylor 展开与 Hessian 定义
3. 最后补 Local minima 与 Saddle point 的对比图
\end{lstlisting}

\begin{figure}[htbp]
  \centering
  \includegraphics[width=0.72\linewidth]{contents/week01/lesson04-local-minima-and-saddle-points/figures/lesson-figure.jpg}
  \caption{临界点章节图片占位}
  \label{fig:w01l04-placeholder}
\end{figure}

\end{document}
